\section{Software}

\subsection{Auslesen einer entwickelten Hardware}
Aufgabe ist es, eine von Christoph Schwaninger vorbereitete Hardware zur Figurenerkennung für die Spiellogik zu entwickeln, dass 
zuverlässig unabhängig voneinander platzierte Figuren auf den Sensoren erkannt werden und eine entsprechende Nachricht 
in der Konsole ausgegeben wird.

\subsection{Rotary Encoder als Eingabegerät}
Ziel ist es, einen Rotary Encoder Code zu schreiben, um sowohl die Drehbewegung als auch den eingebauten Button 
gleichfalls zu erkennen und in der Konsole auszugeben. Außerdem soll erklärt werden, wie ein Rotary Encoder grundlegend 
funktioniert.

\subsection{Entwicklung eines Menüs}
Es soll ein intuitives Menü entwickelt werden, welches nach Einschalten des Schachbretts erscheint. Das Menü soll 
anschließend so mit dem Rotary Encoder verbunden werden, dass man problemlos navigieren kann.

\subsection{Einbindung der Spielererkennung}
Aufgabe ist es, die Spielererkennung über einen RFID-Sensor zu implementieren. Für Testzwecke soll zunächst
eine Nachricht in der Konsole ausgegeben werden, wenn ein NFC-Tag erkannt wird.

\subsubsection{Spiellogik}
Die Spiellogik soll in zwei Teile geteilt werden: Einerseits die Logik für das Einzelspiel gegen einen Schachcomputer, 
für den Stockfish verwendet werden soll. Andererseits die Logik für den Mehrspielermodus gegen eine andere Person. 
Bei Letzterem soll das Spiel aufgezeichnet und in eine PGN-Datei gespeichert werden.

\subsection{Die Datenübertragung}
Nach Beendigung einer Partie im Multiplayer soll die PGN-Datei an einen Server, welcher vorher von Christina Seidner 
vorbereitet wurde, übertragen werden.